\section{Effective asymptotics and decimal digits}
\label{sec:asymptotics}

The expression $u_n\sim v_n$ is not taken as primitive notation.  We prove
rational inequalities with explicit moduli for the Cauchy codes of $\pi$ and
square roots.

\subsection{A Wallis bound and the leading equivalent}

Let
\[
 q_n=\frac1{4^n}\binom{2n}{n}.
\]
The finite Wallis inequalities give, for $n\ge1$,
\[
 \frac1{\sqrt{\pi(n+1/2)}}<q_n<\frac1{\sqrt{\pi n}}.             \tag{8.1}\label{eq:wallis-central}
\]
For completeness, one constructive derivation starts with
$I_m=\int_0^{\pi/2}\sin^m t\,dt$.  Integration by parts gives
$I_m=(m-1)I_{m-2}/m$, while
$I_{2n+1}<I_{2n}<I_{2n-1}$ follows pointwise on $(0,\pi/2)$.  Substitute the
finite product formulas for $I_{2n}$ and $I_{2n+1}$, square, and clear all
positive factors.  Refining the elementary comparison with
$\sin t$ on rational polygonal meshes gives Equation~\eqref{eq:wallis-central}
and a primitive-recursive Cauchy code for $\pi$.  No impredicative completeness
principle enters the proof.

Define the normalized ratio
\[
 R_n=\frac{\sqrt\pi\,n^{3/2}\Cat_n}{4^n}
     =\frac n{n+1}\sqrt{\pi n}\,q_n.                            \tag{8.2}
\]
Equation~\eqref{eq:wallis-central} implies
\[
 \frac n{n+1}\sqrt{\frac n{n+1/2}}<R_n<\frac n{n+1}<1.          \tag{8.3}\label{eq:Rbounds}
\]
The rational inequalities
\[
 \frac n{n+1}\ge1-\frac1n,
 \qquad
 \frac1{\sqrt{1+x}}\ge1-\frac x2\quad(x\ge0)
\]
then give
\[
 0<1-R_n<\frac{5}{4n}<\frac2n.                                  \tag{8.4}\label{eq:leading-error}
\]
The square-root inequality is proved by squaring after checking both sides are
nonnegative.  Therefore a convenient primitive-recursive modulus is
\[
 n\ge2q\quad\Longrightarrow\quad \abs{R_n-1}\le\frac1q.         \tag{8.5}\label{eq:leading-modulus}
\]
This proves, in the exact constructive sense,
\[
 \boxed{\displaystyle
 \Cat_n\sim\frac{4^n}{\sqrt\pi\,n^{3/2}}.}                       \tag{8.6}
\]

\subsection{The first three displayed asymptotic terms}

The source page prints
\[
 \Cat_n\sim\frac{4^n}{\sqrt\pi}
 \left(n^{-3/2}-\frac98n^{-5/2}+\frac{145}{128}n^{-7/2}+\cdots\right).
\tag{8.7}\label{eq:printed-expansion}
\]
We now give an effective remainder.

\begin{lemma}[Two-term effective Stirling logarithm]
For $m\ge1$,
\[
 \log(m!)=
 \left(m+\frac12\right)\log m-m+\frac12\log(2\pi)
 +\frac1{12m}-\frac1{360m^3}+\rho_m,                            \tag{8.8}
\]
where
\[
 0<\rho_m<\frac1{1260m^5}.                                      \tag{8.9}
\]
\end{lemma}

\begin{proof}
Apply Euler--Maclaurin to $\log x$ through the Bernoulli-$B_4$ term.  The next
derivative has fixed sign, so the integral remainder lies between zero and
the first omitted Bernoulli term.  At this fixed order the Bernoulli
polynomials, their maxima on $[0,1]$, and the integral comparison are rational
polynomial calculations.  For an object-language proof, replace the integral
by upper and lower rational Riemann sums; the derivative bound gives a mesh
modulus primitive recursive in $(m,q)$.  The endpoint constant is fixed by the
Wallis code for $\pi$.  This is the standard effective proof of Stirling, but
all quantifiers have explicit finite bounds.
\end{proof}

Subtracting twice the formula at $n$ from the formula at $2n$ gives
\[
 \log\left(\frac{\binom{2n}{n}\sqrt{\pi n}}{4^n}\right)
 =-\frac1{8n}+\frac1{192n^3}+\eta_n,
 \qquad \abs{\eta_n}\le\frac1{600n^5}.                          \tag{8.10}\label{eq:central-log}
\]
Put $u=1/n$.  Since $R_n$ is the central-binomial ratio multiplied by
$1/(1+u)$,
\[
 \log R_n=-\frac98u+\frac12u^2-\frac{21}{64}u^3+\delta_n,
 \qquad \abs{\delta_n}\le\frac{13}{50}u^4                     \tag{8.11}\label{eq:Rlog}
\]
for $n\ge8$.  Here the alternating expansion of $\log(1+u)$ has remainder at
most $u^4/4$, and Equation~\eqref{eq:central-log} supplies the smaller
$u^5/600$ term.

For $n\ge8$, Equation~\eqref{eq:Rlog} gives $\abs{\log R_n}<3/20$.
Here is an explicit rational audit of the Taylor constant.  Put
$z=\log R_n$, $a_1=-9/8$, and $a_2=1/2$.  Since $u\le1/8$,
Equation~\eqref{eq:Rlog} implies
\[
 |z|\le\frac65u,\qquad
 |z-a_1u|\le\frac{11}{20}u^2,\qquad
 |z+a_1u|\le\frac{12}{5}u.
\]
Also $e^{3/20}<7/6$, as follows by summing its positive Taylor series with a
geometric tail after the linear term.  Taylor's theorem therefore gives
\begin{align*}
 \left|e^z-\left(1+a_1u+
          \left(a_2+\frac{a_1^2}{2}\right)u^2\right)\right|
 &\le \left(\frac{21}{64}+\frac{13}{400}\right)u^3
   +\frac12\frac{11}{20}\frac{12}{5}u^3
   +\frac{7}{36}\left(\frac65\right)^3u^3\\
 &<2u^3.
\end{align*}
The three summands respectively bound the cubic-and-$\delta_n$ part of $z$,
the correction from $z^2-a_1^2u^2$, and the exponential remainder.  Hence
\[
 \boxed{\displaystyle
 \abs{R_n-
  \left(1-\frac9{8n}+\frac{145}{128n^2}\right)}
 \le\frac2{n^3}\qquad(n\ge8).}                                 \tag{8.12}\label{eq:three-term-bound}
\]
The coefficient calculation is
\[
 \frac12+\frac12\left(\frac98\right)^2=\frac{145}{128}.
\]
Collecting one further coefficient gives the next term
$-1155/(1024n^3)$.  It is not needed below; the report stops at
Equation~\eqref{eq:three-term-bound} because that bound exactly certifies all
terms printed by MathWorld.

Every transcendental operation in this argument is evaluated by a rational
interval.  For a requested accuracy $1/q$, choose the Taylor lengths by the
explicit alternating/geometric tail bounds, and choose $n$ from the displayed
polynomial inequalities.  Thus the proof has a primitive-recursive modulus
and is admissible in the constructive setting.

\subsection{Decimal digit ratios}

Let $d(N)$ be the number of decimal digits of $\Cat_N$, so
$d(N)=\floor{\log_{10}\Cat_N}+1$.  From Equation~\eqref{eq:Rbounds}, for
$N\ge2$,
\[
 \log_{10}\Cat_N
 =N\log_{10}4-\frac32\log_{10}N-\frac12\log_{10}\pi
  +\log_{10}R_N.                                                 \tag{8.13}
\]
Since $1/2<R_N<1$ for $N\ge3$, the last two terms have bounded absolute value.
Together with the one-unit floor error, a convenient explicit estimate is
\[
 \abs{\frac{d(N)}N-\log_{10}4}
 \le\frac{\tfrac32\log_{10}N+3}{N}\qquad(N\ge3).                \tag{8.14}\label{eq:digit-bound}
\]
In particular, at $N=10^m$ with $m\ge1$,
\[
 \abs{\frac{d(10^m)}{10^m}-\log_{10}4}
 \le\frac{\tfrac32m+3}{10^m},                                  \tag{8.15}\label{eq:digit-ratio}
\]
which proves the convergence claimed by MathWorld with a primitive-recursive
modulus.  The logarithm base ten is a rational Cauchy code obtained from the
series for $\log(1+x)$ after range reduction.  The finite digit values printed
on the source page are computational data; the supplied checker independently
certifies those within its configured feasible range, while
Equation~\eqref{eq:digit-ratio} proves the general limiting statement.
